% 中文摘要

随着云计算、微服务和容器化技术的广泛应用，现代软件系统规模日益庞大、架构日趋复杂，传统运维方法依赖人工经验，难以高效应对海量日志、监控指标和告警信息。为提升故障诊断效率与运维智能化水平，本文设计并实现了一个基于大语言模型的智能运维故障诊断原型系统——LLM-Ops Copilot。

本文首先构建了面向运维场景的领域知识库。基于ChromaDB向量数据库，构建了包含标准操作流程（SOP）文档、历史故障案例和系统架构文档的运维知识库，设计了层级递归切分策略和混合检索方法，实现了文档的向量化存储与语义检索，支持按服务名称和故障类型进行元数据过滤，增强了模型获取任务相关知识的能力。其次，提出了多源运维数据融合方法。设计了证据融合器（EvidenceBuilder）模块，包含三级日志处理流水线、异常指标识别与描述、关联分析和严重程度评估等子步骤，将知识检索结果、日志摘要和监控指标描述统一组织为结构化的分析上下文。再次，设计并实现了基于大模型的根因分析引擎（RCA Engine）。设计了包含角色设定、上下文规范、推理约束和输出格式的完整结构化Prompt模板，引导模型按照"故障摘要—证据分析—根因假设—修复建议"的逻辑顺序展开分析，支持置信度评分、多根因假设和数据缺口识别等功能。最后，通过7个典型故障场景和四组对照实验验证了方法的有效性。

实验结果表明：仅依赖大语言模型直接回答时，系统在专业性方面存在明显不足，准确率仅为43\%；引入RAG知识增强后提升至57\%；进一步融合实时日志和监控指标后提升至71\%；完整方案达到86\%，Top-3准确率达到100\%，证据覆盖率达到78\%。典型案例分析验证了证据融合和结构化Prompt对于提升诊断质量的关键作用，为智能运维提供了新的技术路径。
